\documentclass[11pt]{article}

% ---------------------------------------------------------------------------
% Packages
% ---------------------------------------------------------------------------
\usepackage[margin=1in]{geometry}
\usepackage{amsmath,amssymb,amsthm}
\usepackage{bm}
\usepackage{booktabs}
\usepackage{algorithm}
\usepackage{algpseudocode}
\usepackage{xcolor}
\usepackage{graphicx}
\usepackage{enumitem}
\usepackage[hidelinks]{hyperref}

% ---------------------------------------------------------------------------
% Notation macros
% ---------------------------------------------------------------------------
\newcommand{\R}{\mathbb{R}}
\newcommand{\Ehat}{\hat{\mathbb{E}}}
\newcommand{\obs}{s}
\newcommand{\act}{a}
\newcommand{\parMLP}{\Phi}                 % ParallelMLP feature bank
\newcommand{\mask}{\bm{m}}                 % k-hot mask
\newcommand{\masksoft}{\tilde{\bm{m}}}     % STE (soft) mask
\newcommand{\Ho}{H}                        % number of parallel sub-networks
\newcommand{\kk}{k}                        % active sub-networks
\newcommand{\kwta}{\operatorname{KWTA}}
\newcommand{\topk}{\operatorname{top\text{-}k}}
\newcommand{\code}[1]{\texttt{\small #1}}

\newtheorem{definition}{Definition}

\title{\bfseries Adapting Functional Task Networks to\\ Offline Reinforcement Learning}
\author{Kevin L.\ McKee\\ \normalsize Astera Institute}
\date{\today}

\begin{document}
\maketitle

\begin{abstract}
Functional Task Networks (FTN) are a continual-learning architecture in which a
single shared bank of small parallel sub-networks is specialized to a task not by
retraining its weights but by selecting a sparse, per-task binary \emph{mask} over
the sub-networks that feed a linear readout. Tasks that share structure share
active sub-networks; tasks that do not are routed through disjoint sub-networks,
which is what makes the architecture resistant to catastrophic forgetting. This
document is the technical reference for three concrete adaptations of FTN to
value-based and policy-based learning from stored (off-policy / offline)
experience: (i) \textbf{FTN--DQN}, a discrete-action deep $Q$-network; (ii)
\textbf{FTN--DDPG}, a continuous-control deterministic actor--critic; and (iii)
\textbf{FTN--BC}, offline behavior cloning of expert demonstrations. All three
share the same FTN backbone, the same annealed straight-through mask search, and
the same mask-overlap-filtered cross-task replay. We describe exactly how each
piece is implemented---the network modules, the mask-reconfiguration objectives,
the replay/consolidation machinery, and the per-phase training loops---so that
the three implementations can be read, reproduced, and extended. Results and
benchmark tasks are deliberately out of scope; this is documentation of
\emph{method}, not of performance.
\end{abstract}

\tableofcontents

% ===========================================================================
\section{Overview}
% ===========================================================================

The three techniques documented here all instantiate the same idea in three
different learning regimes. A shared, fixed-weight backbone---a bank of $\Ho$
small independent multilayer perceptrons whose scalar outputs are gated by a
$\kk$-hot binary mask before a linear readout---is specialized to each new task by
\emph{searching for a mask}, not by fine-tuning weights. Learning a sequence of
tasks $1,2,\dots,N$ therefore decomposes, per task, into two operations:

\begin{enumerate}[leftmargin=2em,itemsep=2pt]
  \item \textbf{Mask reconfiguration.} Given a small support set of experience for
  the current task, find a $\kk$-hot mask $\mask \in \{0,1\}^{\Ho}$,
  $\bm{1}^\top\mask = \kk$, that minimizes a task-appropriate loss. The mask is
  found by a stochastic, annealed straight-through estimator (STE) search
  (\S\ref{sec:ste}), which optimizes continuous logits while committing to a hard
  top-$k$ selection at every step.
  \item \textbf{Backbone update under the committed mask.} With the mask fixed,
  take gradient steps on the backbone weights using the learning rule of the
  regime---TD for DQN, deterministic policy gradient for DDPG, or supervised MSE
  for BC. Crucially, gradients flow only through the active sub-networks, so a
  task can only disturb the sub-networks its mask selects.
\end{enumerate}

Cross-task retention is handled by a single shared mechanism
(\S\ref{sec:consolidation}): archived experience from past phases is stored
\emph{tagged with the mask it was collected under}, and during the current phase
we additionally replay past chunks whose mask \emph{partially overlaps} the
current one. Because those chunks are forwarded under their own masks, the
resulting gradients land on exactly the sub-networks shared between old and new
tasks, consolidating them, while leaving disjoint sub-networks untouched.

Table~\ref{tab:map} maps each technique to its source file and summarizes what
changes across the three. The remainder of the document is organized as:
\S\ref{sec:backbone} the shared backbone; \S\ref{sec:ste} the mask search;
\S\ref{sec:consolidation} the shared consolidation mechanism and the ablation
baselines; \S\ref{sec:dqn} FTN--DQN; \S\ref{sec:ddpg} FTN--DDPG; \S\ref{sec:bc}
FTN--BC.

\begin{table}[t]
\centering
\small
\begin{tabular}{@{}llll@{}}
\toprule
 & \textbf{FTN--DQN} & \textbf{FTN--DDPG} & \textbf{FTN--BC}\\
\midrule
Source file & \code{train\_grid\_dqn.py} & \code{train\_mt10.py} & \code{train\_mt10\_expert.py}\\
            &                            & \code{train\_cw.py}    & \\
Networks    & 1 $Q$-net                  & actor + critic (+targets) & actor only\\
Masks / task& 1 ($\mask$)                & 2 ($\mask_a,\mask_c$)     & 1 ($\mask_a$)\\
Learning rule& TD($Q$)                   & deterministic PG          & supervised MSE\\
Data source & online replay ring          & online replay ring        & offline expert demos\\
Reconfig obj.& $Q$ vs.\ target/returns    & critic vs.\ returns, then $-Q$ & BC MSE\\
Consolidation& per-phase replay, one mask & per-chunk replay, mask pair & per-task support, one mask\\
\bottomrule
\end{tabular}
\caption{The three FTN adaptations and where they differ. All share the backbone
of \S\ref{sec:backbone}, the mask search of \S\ref{sec:ste}, and the
overlap-filtered consolidation of \S\ref{sec:consolidation}.}
\label{tab:map}
\end{table}

% ===========================================================================
\section{The shared FTN backbone}
\label{sec:backbone}
% ===========================================================================

All three techniques are built on the modules in \code{ftn.py}. The backbone has
two components: a bank of parallel sub-networks that maps an input to $\Ho$ scalar
features, and a masked linear readout.

\subsection{ParallelMLP: a bank of $\Ho$ independent sub-networks}

The class \code{ParallelMLP} evaluates $\Ho$ small MLPs in parallel, each mapping
the same input vector $x \in \R^{d_{\mathrm{in}}}$ to a scalar. Every sub-network
has \code{layers} hidden layers of width \code{inner\_size} and a scalar output,
so the internal size schedule is
\[
\texttt{sizes} = [\,d_{\mathrm{in}},\ \underbrace{d_h,\dots,d_h}_{\text{layers}},\ 1\,],
\qquad d_h = \texttt{inner\_size}.
\]
The weights are stored as a stack of 3-D parameter tensors, one per layer
transition $\ell$, of shape $(\Ho,\ \texttt{sizes}[\ell],\ \texttt{sizes}[\ell{+}1])$,
with a matching $(\Ho,\ \texttt{sizes}[\ell{+}1])$ bias. The forward pass batches
over both the data dimension and the $\Ho$ sub-networks with a single einsum per
layer:
\[
h^{(\ell+1)}_{b,\,i,\,j} \;=\; \sum_{p} h^{(\ell)}_{b,\,i,\,p}\, W^{(\ell)}_{i,\,p,\,j} \;+\; b^{(\ell)}_{i,\,j},
\]
i.e.\ \code{einsum("bki,kij->bkj", h, W)}, where $b$ indexes the batch and
$i \in \{1,\dots,\Ho\}$ indexes the sub-network. ReLU is applied after every layer
except the last, and the trailing singleton output dimension is squeezed, giving a
feature bank
\[
\parMLP(x) \;=\; z \;\in\; \R^{B \times \Ho}.
\]
Weights use a custom uniform initialization scaled by fan-in: for hidden layers
the bound is $\sqrt{6/\text{fan\_in}}$ and for the output layer
$\sqrt{6/(\text{fan\_in}+\text{fan\_out})}$; biases are drawn from
$\mathcal{U}(\pm 0.1/\sqrt{\text{fan\_in}})$. Because the $\Ho$ sub-networks never
mix, a mask that zeroes a sub-network's feature removes that sub-network's
influence on the readout completely; only the readout weights and the masked
sub-networks receive gradient.

\subsection{Masked linear readout}
\label{sec:readout}

Given the feature bank $z \in \R^{B\times\Ho}$ and a mask $\mask\in\{0,1\}^{\Ho}$,
the readout applies dropout, gates the features element-wise by the (broadcast)
mask, and applies a bias-free linear map:
\[
y \;=\; W_{\mathrm{out}}\,\bigl(\operatorname{dropout}(z)\,\otimes\, \mask\bigr)^{\!\top}.
\]
Dropout is a train-time regularizer only; at rollout/eval the network is put in
eval mode so action selection is deterministic. The three techniques differ only
in the shape of $W_{\mathrm{out}}$ and the nonlinearity that follows:
\begin{itemize}[leftmargin=1.5em,itemsep=1pt]
  \item \textbf{\code{FTNActor}} ($W_{\mathrm{out}}\!:\Ho\to d_a$): output is
  $\tanh(y)$, a bounded continuous action; input is the observation $\obs$.
  \item \textbf{\code{FTNCritic}} ($W_{\mathrm{out}}\!:\Ho\to 1$): input is the
  concatenation $[\obs,\act]$; output is a scalar $Q$-value.
  \item \textbf{\code{FTNQNetwork}} ($W_{\mathrm{out}}\!:\Ho\to |\mathcal{A}|$):
  input is $\obs$; output is a vector of $Q$-values, one per discrete action.
\end{itemize}
Because the mask multiplies the features (not the weights), the same set of
backbone weights realizes different functions for different masks: the mask
\emph{routes} the computation. The soft mask used during search
(\S\ref{sec:ste}) enters at exactly the same place, so the search objective is
differentiable with respect to it.

The $Q$-network additionally provides a batched forward, \code{forward\_batched},
that evaluates $E$ masks at once against a shared support batch: the expensive
$\parMLP$ forward and dropout run \emph{once}, and only the per-mask gate and
final linear are duplicated, producing an $(E, B, |\mathcal{A}|)$ tensor. This is
what makes the batched mask search of \S\ref{sec:ste} cheap.

% ===========================================================================
\section{Annealed straight-through mask search}
\label{sec:ste}
% ===========================================================================

Mask reconfiguration is the operation that specializes the shared backbone to a
task. It is implemented once, in \code{anneal\_ste\_mask}, and reused by all three
techniques with only the loss functional changing.

\subsection{Problem and estimator}

We seek a $\kk$-hot mask minimizing a task loss $\mathcal{L}$ that is only defined
through a forward pass of a masked network:
\[
\mask^\star \;=\; \arg\min_{\mask\in\{0,1\}^{\Ho},\ \bm 1^\top\mask=\kk}\ \mathcal{L}(\mask).
\]
This is a combinatorial selection over $\binom{\Ho}{\kk}$ subsets. We relax it with
a \emph{straight-through} estimator over a vector of continuous logits
$\bm\theta\in\R^{\Ho}$. At each step we form a hard top-$k$ selection and pass
gradients through as if it were the identity:
\begin{align}
\mask_{\mathrm{hard}} &= \kwta\!\bigl(\bm\theta + \bm\epsilon,\ \kk\bigr), \qquad \bm\epsilon\sim\mathcal N(\bm 0,\ \tau^2 \bm I),\\
\masksoft &= \mask_{\mathrm{hard}} + \bm\theta - \operatorname{sg}(\bm\theta),
\label{eq:ste}
\end{align}
where $\kwta(\cdot,\kk)$ returns the binary indicator of the top-$\kk$ entries and
$\operatorname{sg}$ is stop-gradient. In the forward pass $\masksoft =
\mask_{\mathrm{hard}}$ (a hard mask); in the backward pass
$\partial\masksoft/\partial\bm\theta = \bm I$, so $\nabla_{\bm\theta}\mathcal{L}$
is the gradient of the loss with respect to the (gated) features. The network is
therefore always evaluated under a genuine $\kk$-hot mask, but the logits still
receive a smooth learning signal.

\subsection{Annealing schedule}

The noise $\bm\epsilon$ added \emph{before} the top-$k$ is what lets the search
explore alternative subsets early and settle late. Its scale is annealed
geometrically. With \code{n\_steps} total steps and $\text{denom}=\max(1,
\texttt{n\_steps}/4)$, at step $t$
\[
\tau_t \;=\; \sigma_t \cdot \exp\!\left(-\,t/\text{denom}\right),\qquad
\sigma_t = \max\bigl(\operatorname{std}(\bm\theta),\ 0.1\bigr),
\]
so the perturbation is scaled to the current spread of the logits and decays to a
near-deterministic top-$k$ by the end of the search. The logits are optimized by
plain SGD (\code{reconfig\_lr}). Throughout, the routine tracks the
\emph{lowest-loss hard mask seen} and returns it, together with the final logits
(for warm-starting), the initial loss, and the final loss. Warm-starting
initializes $\bm\theta$ from a previous mask's logits scaled by a small noise
factor; cold-starting initializes from $\texttt{noise\_scale}\cdot\mathcal
N(\bm 0,\bm I)$. Algorithm~\ref{alg:ste} gives the exact procedure.

\begin{algorithm}[t]
\caption{\code{anneal\_ste\_mask}: annealed STE search for one $\kk$-hot mask}
\label{alg:ste}
\begin{algorithmic}[1]
\Require sizes $\Ho,\kk$; loss $\mathcal L(\cdot)$; steps $T$; lr $\eta$; init logits $\bm\theta_0$ or \textsc{None}; noise scale $c$
\If{$\bm\theta_0$ is \textsc{None}} $\bm\theta \gets c\cdot\mathcal N(\bm 0,\bm I_{\Ho})$ \Comment{cold start}
\Else\ $\bm\theta \gets c\cdot\bm\theta_0$ \Comment{warm start} \EndIf
\State $\text{best}\gets\varnothing,\ \text{lowest}\gets+\infty,\ \text{denom}\gets\max(1, T/4)$
\For{$t = 0,\dots,T-1$}
  \State $\sigma \gets \max(\operatorname{std}(\bm\theta),\ 0.1)$;\quad $\tau \gets \sigma\,e^{-t/\text{denom}}$
  \State $\mask_{\mathrm{hard}} \gets \kwta(\bm\theta + \tau\cdot\mathcal N(\bm 0,\bm I),\ \kk)$
  \State $\masksoft \gets \mask_{\mathrm{hard}} + \bm\theta - \operatorname{sg}(\bm\theta)$ \Comment{Eq.~\eqref{eq:ste}}
  \State $L \gets \mathcal L(\masksoft)$
  \If{$L < \text{lowest}$} $\text{lowest}\gets L$; $\text{best}\gets \mask_{\mathrm{hard}}$ \EndIf
  \State SGD step on $\bm\theta$ with $\nabla_{\bm\theta} L$
\EndFor
\State \Return best hard mask, final logits $\bm\theta$, initial loss, lowest loss
\end{algorithmic}
\end{algorithm}

\subsection{Batched search}
\label{sec:ste-batched}

For eval-time recovery we often want several \emph{independent} cold-started
searches on the same support (to average over the stochasticity of the search).
\code{anneal\_ste\_mask\_batched} runs $E$ searches at once: the logits are an
$(E,\Ho)$ matrix, the loss functional returns a per-row $(E,)$ vector, and the sum
of the per-row losses is back-propagated in a single graph while each row tracks
its own best hard mask. Combined with \code{forward\_batched}
(\S\ref{sec:readout}), the shared support is forwarded through $\parMLP$ once for
all $E$ rows. FTN--DQN uses this at eval; FTN--DDPG and FTN--BC call the scalar
search per eval episode instead.

\subsection{Lateral spatial smoothing (optional)}
\label{sec:smooth}

FTN--BC additionally post-processes each searched mask with the
\code{LateralSmoother} (a port of \emph{SmoothKWTA}). The flat $\kk$-hot mask is
laid onto a square grid of $\texttt{square\_dim}\times\texttt{square\_dim}$ blocks
(requiring $\Ho$ to be a multiple of $\texttt{square\_dim}^2$), smoothed by
\code{lateral\_steps} passes of a uniform-kernel 2-D convolution mixed back into
the mask by a Polyak coefficient, and re-sparsified with top-$\kk$:
\[
\mask \gets \kwta\!\Bigl(\underbrace{\lambda\,\mask + (1-\lambda)\,\operatorname{conv}(\mask)}_{\text{repeated \texttt{lateral\_steps} times}},\ \kk\Bigr).
\]
This clusters a task's active sub-networks into a contiguous spatial blob rather
than a random scatter, which encourages neighboring tasks to reuse neighboring
sub-networks. Smoothing is a no-op (identity) when
$\texttt{lateral\_steps}\le 0$ or $\texttt{lateral\_polyak}\ge 1$; it is only
wired into FTN--BC.

% ===========================================================================
\section{Cross-task consolidation and baselines}
\label{sec:consolidation}
% ===========================================================================

The mechanism that lets a task be learned without destroying earlier tasks is
shared across all three techniques, differing only in bookkeeping granularity.

\subsection{Mask-overlap-filtered replay}

Every completed phase archives its experience tagged with the mask(s) it was
collected under. When learning the current task with mask $\mask_{\mathrm{cur}}$,
we compute the overlap of each archived mask $\mask_{\mathrm{old}}$,
\begin{equation}
\rho(\mask_{\mathrm{cur}},\mask_{\mathrm{old}}) \;=\; \frac{\mask_{\mathrm{cur}}^\top \mask_{\mathrm{old}}}{\max(\bm 1^\top\mask_{\mathrm{cur}},\,1)}
\;=\; \frac{|\,\text{active}_{\mathrm{cur}}\cap\text{active}_{\mathrm{old}}\,|}{\kk},
\label{eq:rho}
\end{equation}
(the last equality holds for fixed-$\kk$ masks), and \emph{replay only those
archived chunks whose overlap lies in the band $(\rho_{\mathrm{low}},
\rho_{\mathrm{high}}]$}. The rationale is a three-way split:
\begin{itemize}[leftmargin=1.5em,itemsep=1pt]
  \item $\rho \le \rho_{\mathrm{low}}$: the old task uses essentially disjoint
  sub-networks; the current update cannot disturb it, so no replay is needed.
  \item $\rho > \rho_{\mathrm{high}}$ (in particular $\rho=1$, the same mask):
  the data is effectively current-task data and already covered by the ordinary
  replay ring; replaying it adds nothing.
  \item $\rho \in (\rho_{\mathrm{low}}, \rho_{\mathrm{high}}]$: the tasks
  \emph{partially} share sub-networks. Here a current-task gradient step could
  degrade the old task through the shared sub-networks, so we replay the old
  chunk. Because we forward it \emph{under its own mask}, its gradient lands on
  exactly the shared sub-networks and pins their behavior to what the old task
  needs, consolidating the overlap.
\end{itemize}
The crucial implementation detail is that archived data is always forwarded under
the mask it belongs to, never under the current mask---so the data-to-mask
correspondence is exact and the gradient is a legitimate old-task gradient.

\subsection{Ablation baselines}
\label{sec:baselines}

FTN--DDPG and FTN--BC implement three reference conditions that share the entire
pipeline except for how masks are chosen, so result cells are directly
comparable:
\begin{description}[leftmargin=1.5em,itemsep=1pt]
  \item[\code{no\_mask}] An all-ones mask ($\kk=\Ho$), frozen. The readout sees
  every sub-network; this is plain dense learning over the backbone and is the
  catastrophic-forgetting reference. Reconfig and consolidation are disabled
  (identical masks give $\rho=1$, outside the band).
  \item[\code{fixed\_mask}] An oracle disjoint allocation: task $i$ owns
  sub-networks $[\,i\kk,(i{+}1)\kk)$, requiring $N\kk\le\Ho$, recalled exactly at
  eval by (oracle) task identity. Disjoint masks give $\rho=0$, so there is no
  consolidation and forward-transfer eval cells are undefined.
  \item[\code{stored\_mask}] Full FTN training, but at eval the mask is recalled
  from what the task used during its own training phase rather than re-derived
  from support. This isolates mask \emph{recovery} from mask \emph{storage}.
\end{description}
The default (full FTN) condition instead \emph{recovers} each task's mask at eval
by a fresh cold-started search on freshly gathered support, so it needs no oracle
task identity.

% ===========================================================================
\section{FTN--DQN: discrete-action value learning}
\label{sec:dqn}
% ===========================================================================

FTN--DQN (\code{train\_grid\_dqn.py}) applies FTN to a discrete-action deep
$Q$-network. There is a single network, \code{FTNQNetwork}, and a single mask
$\mask$ per task; the target network is a Polyak-averaged copy.

\subsection{TD update under a fixed mask}

With mask $\mask$ committed, one gradient step is the standard DQN TD update, but
every forward pass is masked. For a minibatch of transitions $(\obs,\act,r,\obs')$
with terminal flag $d$,
\begin{align}
y &= r + \gamma\,(1-d)\,\max_{\act'} Q_{\bar\phi}(\obs',\act';\,\mask),\\
\mathcal L_{\mathrm{TD}} &= \Ehat\Bigl[\bigl(Q_\phi(\obs,\act;\,\mask) - y\bigr)^2\Bigr],
\end{align}
where $Q_{\bar\phi}$ is the target network (same mask). We take an Adam step on
$\phi$, then Polyak-update the target: $\bar\phi \gets (1-\tau)\bar\phi +
\tau\phi$. The masked $\max$ and $\operatorname{gather}$ operate on the
$|\mathcal A|$-vector output of the readout. (The loss scalar is not synced back
to the host, since this update runs on the order of $10^5$ times per sweep.)

\subsection{Mask reconfiguration objectives}
\label{sec:dqn-recon}

The reconfig loss $\mathcal L(\masksoft)$ passed to Algorithm~\ref{alg:ste} scores
a candidate mask by how well the masked $Q$-net explains a rewarded support set.
The support provides $(\obs,\act,r,\obs',d)$ and Monte-Carlo returns $G_t$
(computed on the concatenated support, respecting episode boundaries via $d$).
Let $q_{\obs} = \max_{\act'} Q_\phi(\obs,\act';\masksoft)$ be the masked greedy
value. \code{reconfig\_objective} selects one of:
\[
\begin{array}{ll}
\text{\code{td}} & \Ehat\bigl[(q_\obs - (r+\gamma(1-d)\max_{\act'}Q_{\bar\phi}(\obs',\act';\masksoft)))^2\bigr]\\[2pt]
\text{\code{mse\_returns}} & \Ehat\bigl[(q_\obs - G)^2\bigr]\\[2pt]
\text{\code{cos\_returns}} & -\cos\bigl(q_\obs,\ G\bigr)\quad\text{(scale-invariant)}\\[2pt]
\text{\code{r\_returns}} & -\operatorname{Pearson}\bigl(q_\obs,\ G\bigr)\quad\text{(scale- \& offset-invariant; default)}
\end{array}
\]
The default \code{r\_returns} is the negative Pearson correlation---i.e.\ a
mean-centered cosine similarity---so the search need only find a mask under which
the $Q$-values \emph{rank states the way the returns do}, without matching their
magnitude or offset. This is robust to the arbitrary scale of a freshly-masked,
untrained readout. The bootstrap (\code{td}) target is computed under the same
soft mask inside \code{no\_grad}, so gradients flow only through the prediction.

\subsection{Per-phase training loop}
\label{sec:dqn-loop}

For each task in sequence (\code{train\_grid\_dqn}):
\begin{enumerate}[leftmargin=2em,itemsep=1pt]
  \item \textbf{Exploration.} Roll out up to \code{explore\_episodes} episodes,
  stopping after \code{n\_successful\_explore} rewarded ones. The exploration
  policy is uniform random at task 0, or (\code{warmup\_policy}) either random or
  the previous task's greedy mask thereafter.
  \item \textbf{Boundary reconfig.} If any reward was found, run
  Algorithm~\ref{alg:ste} on the rewarded support with the chosen objective and
  \emph{commit} the mask; otherwise keep the prior mask. Push the exploration
  episodes into the replay ring.
  \item \textbf{Exploitation.} For \code{n\_trials\_per\_task} episodes: collect
  one $\varepsilon$-greedy episode under the current mask, push it, and (once the
  ring has $\ge$ \code{batch\_size} transitions) run \code{updates\_per\_step} TD
  updates per env step. After each ordinary update, run up to
  \code{mask\_replay\_n\_per\_step} \emph{cross-task} updates: draw a filtered
  past phase (\S\ref{sec:dqn-replay}) and take a TD step \emph{under that phase's
  mask}.
  \item \textbf{Mid-phase reconfig.} Every \code{reconfig\_every\_trials} trials,
  re-search and commit the mask on either the most recent episodes
  (\code{reconfig\_support="recent"}) or a fresh random-exploration support.
  \item \textbf{Archive \& eval.} Store the phase's mask, archive its ring
  contents tagged with that mask, then eval every task (\S\ref{sec:dqn-eval}).
\end{enumerate}

\subsection{Single-mask overlap-filtered replay}
\label{sec:dqn-replay}

\code{MaskedReplayBufferDQN} is a GPU-resident ring for the current phase plus a
per-phase archive. Discrete actions are stored as \code{int64}; each archived
phase keeps its transitions and the single mask $\mask$ it was collected under, and
a stacked $(P,\Ho)$ matrix of past masks is cached for vectorized overlap.
\code{sample\_filtered} computes Eq.~\eqref{eq:rho} against all past masks in one
op, keeps the phases whose $\rho\in(\rho_{\mathrm{low}},\rho_{\mathrm{high}}]$,
picks \emph{one} qualifying phase uniformly (so the $(\text{batch},
\mask_{\mathrm{old}})$ pair is unambiguous), and samples a batch from it. The
returned mask is fed directly to the TD update (\S\ref{sec:dqn-loop}), so the
old-task gradient only touches sub-networks active under $\mask_{\mathrm{old}}$.

\subsection{Eval: recovering a mask from scratch}
\label{sec:dqn-eval}

Eval never reuses a stored mask (in the default condition). For each task it
gathers a rewarded support---either by random/warmup rollouts, or by the
\code{bootstrap} method, in which each exploration episode after the first reward
acts under a fresh cold-start reconfig on the cumulative rewarded support (so a
bad early mask cannot trap exploration). It then runs \code{eval\_episodes}
independent cold-started searches \emph{batched} into one STE pass
(\S\ref{sec:ste-batched}), rolls one greedy episode under each recovered mask, and
reports the mean success rate. Warm-start logits are cleared for the whole eval so
no recovery is contaminated by the last-trained task, and restored afterward.

% ===========================================================================
\section{FTN--DDPG: continuous-control actor--critic}
\label{sec:ddpg}
% ===========================================================================

FTN--DDPG applies FTN to deterministic-policy-gradient continuous control. It is
implemented twice with an identical agent: \code{train\_mt10.py} (Meta-World MT10,
cheat-initialized resets, random exploration) and \code{train\_cw.py} (the
Continual World CW10/CW20 protocol, scripted-expert support and demonstration
seeding). The differences between the two files are in \emph{environment setup and
support collection}, not in the FTN mechanics; we document the shared agent and
note the CW-specific additions at the end.

\subsection{Two networks, two masks}

The agent (\code{FTNDDPG}) holds an \code{FTNActor} $\mu_\phi$, an \code{FTNCritic}
$Q_\psi$, and Polyak target copies $\mu_{\bar\phi}, Q_{\bar\psi}$. Each task has
\emph{two} masks: an actor mask $\mask_a$ (gating the policy's sub-networks) and a
critic mask $\mask_c$ (gating the value network's). The greedy action is
$\mu_\phi(\obs;\mask_a)$ with the actor in eval mode; exploration adds a noise
sample (\S\ref{sec:ddpg-noise}) before clipping to $[-1,1]$.

\subsection{Two-stage mask reconfiguration}
\label{sec:ddpg-recon}

Reconfiguration searches the two masks in sequence
(\code{reconfig\_masks}):
\begin{enumerate}[leftmargin=2em,itemsep=1pt]
  \item \textbf{Critic mask} $\mask_c$ via the same
  \code{reconfig\_objective} family as FTN--DQN (\S\ref{sec:dqn-recon}), now with
  $q_{\obs} = Q_\psi(\obs,\act;\masksoft)$ scored against MC returns $G_t$
  (default \code{r\_returns}, i.e.\ $-\operatorname{Pearson}(Q,G)$), or a
  DDPG TD target $r+\gamma(1-d)Q_{\bar\psi}(\obs',\mu_{\bar\phi}(\obs'))$ under
  \code{td}.
  \item \textbf{Actor mask} $\mask_a$, holding the freshly-chosen critic mask
  fixed, by maximizing the critic's estimate of the masked policy:
  \[
  \mathcal L_{\mathrm{actor}}(\masksoft) = -\,\Ehat\bigl[\,Q_\psi\bigl(\obs,\ \mu_\phi(\obs;\masksoft);\ \mask_c\bigr)\bigr].
  \]
\end{enumerate}
Both are Algorithm~\ref{alg:ste} searches; both may warm-start from the running
logits. Committing sets $\mask_a,\mask_c$ and stores their logits.

\subsection{Joint update with summed cross-task loss}
\label{sec:ddpg-update}

The core learning step is \code{update\_joint}, which takes the current batch
(forwarded under $\mask_a,\mask_c$) plus a list of \emph{recalled} chunks (each a
$(\text{batch},\mask_a^{(j)},\mask_c^{(j)})$ triple) and optimizes the
\emph{summed} loss over all terms. For the critic, each term contributes its own
TD loss under its own mask:
\[
\mathcal L_{\mathrm{critic}} = \sum_{j}\ \Ehat_j\!\Bigl[\bigl(Q_\psi(\obs,\act;\mask_c^{(j)}) - (r + \gamma(1-d)\,Q_{\bar\psi}(\obs',\mu_{\bar\phi}(\obs';\mask_a^{(j)});\mask_c^{(j)}))\bigr)^2\Bigr],
\]
and for the actor, each term contributes its own $-Q$:
\[
\mathcal L_{\mathrm{actor}} = \sum_{j}\ -\,\Ehat_j\!\bigl[Q_\psi(\obs,\ \mu_\phi(\obs;\mask_a^{(j)});\ \mask_c^{(j)})\bigr].
\]
Every term gets its own forward pass so the mask-to-data match is exact;
gradients accumulate on sub-networks shared across the masks. One critic step, one
actor step, and one Polyak target soft-update per call. (The summation points are
also the natural hook for per-chunk weighting, left as an unweighted sum.)

\subsection{Decoupled collect/train loop}
\label{sec:ddpg-loop}

Exploitation uses an actor--learner-style decoupled loop rather than per-step
interleaving. For \code{n\_iters\_per\_task} iterations:
\begin{enumerate}[leftmargin=2em,itemsep=1pt]
  \item \textbf{Collect.} Step \code{n\_parallel\_train} continuing
  (auto-resetting) envs for \code{collect\_steps} each under the noisy actor, and
  push the burst into the replay ring. Auto-reset fires on
  success/termination/truncation/timeout. A subtle but important distinction is
  made between two flags: the \emph{bootstrap-cutting} $d$ (true MDP terminal or
  absorbing goal only---\emph{not} horizon timeouts, so bootstrapping continues
  through time-limit truncations and values are not biased down near the horizon)
  and an episode-\emph{boundary} flag (every auto-reset) used only to segment MC
  returns so credit never leaks across a reset or across parallel envs.
  \item \textbf{Train.} Run \code{train\_steps} \code{update\_joint} calls. Each
  draws a current batch and, unless in a no-reconfig baseline, up to
  \code{replay\_max\_chunks} recalled chunks via \code{sample\_filtered\_multi}
  (\S\ref{sec:ddpg-replay}); demonstration chunks (\S\ref{sec:ddpg-demo}) are
  summed in on top.
  \item \textbf{Mid-phase reconfig.} Every \code{reconfig\_every\_iters}
  iterations, build support from the latest collect burst (MC returns per env),
  \emph{archive the current segment as a chunk under the outgoing masks}, clear
  the ring, and commit freshly searched masks. This per-interval chunking is what
  guarantees each archived chunk's transitions correspond to exactly one mask
  pair.
\end{enumerate}
Adam state is reset at every task boundary (\code{reset\_optimizer\_state}): stale
momentum must not keep stepping the previous task's sub-networks after they leave
the active mask and their gradients become zero.

\subsection{Per-chunk overlap-filtered replay}
\label{sec:ddpg-replay}

\code{MaskedReplayBufferDDPG} archives \emph{sub-phase chunks}, each tagged with
the $(\mask_a,\mask_c)$ pair it was collected under (the caller archives and clears
at every reconfig commit, so this correspondence is exact).
\code{sample\_filtered\_multi} vectorizes Eq.~\eqref{eq:rho} over all chunks for
both the actor and critic masks and keeps a chunk if \emph{either} its actor
\emph{or} its critic overlap lies in the band; if more than
\code{replay\_max\_chunks} qualify it takes a uniform random subset (a weighting
would bias selection). It returns one $(\text{batch},\mask_a,\mask_c)$ triple per
kept chunk---exactly the input \code{update\_joint} sums over.

\subsection{Demonstration seeding (DDPGfD-style)}
\label{sec:ddpg-demo}

Separately from the $\rho$-filtered archive, the buffer keeps a \emph{protected
demo store}. When \code{demo\_seed} is on, the exploration episodes (which were
collected before the phase's mask existed and so are never pushed to the per-mask
ring) are archived as a demo chunk tagged with the just-committed mask, and
\emph{summed into every joint update} for the phase. Demos are \emph{not}
$\rho$-gated: same-mask data has $\rho=1$ and would be filtered out, but here we
want the agent to learn directly from the seed data under its own mask (e.g.\ from
scripted-expert trajectories). Optionally demos persist across phases
(\code{demo\_seed\_persist}), each replayed under its own task's mask as a
cross-task retention aid.

\subsection{Exploration noise and support collection}
\label{sec:ddpg-noise}

Exploration uses either temporally-correlated Ornstein--Uhlenbeck (OU) noise
(default) or i.i.d.\ Gaussian noise. OU noise,
$dx = \theta(\mu-x)\,dt + \sigma\sqrt{dt}\,\mathcal N(0,I)$, integrates into
sustained directed motion, so as both an additive term on the policy and a
no-policy cold-start driver it actually translates the arm across the workspace
rather than jittering in place; each env owns its own process, reset at episode
boundaries. \code{train\_cw.py} additionally supports \code{scripted} warmup:
per-task MetaWorld expert policies (plus noise) gather rewarded support on hard
manipulation tasks that random exploration never solves---seeding \emph{support
only}; the agent still learns and is evaluated under its own masked policy.

\subsection{Eval}
FTN--DDPG eval mirrors FTN--DQN's: gather rewarded support (random, scripted, or
bootstrap), then for each of \code{eval\_episodes} run a fresh cold-started
reconfig (\code{commit=False}) on that support and one greedy rollout, reporting
mean success. Recall baselines (\code{fixed\_mask}/\code{stored\_mask}) skip
exploration and roll out under the recalled mask; \code{no\_mask} rolls out under
the frozen all-ones mask.

% ===========================================================================
\section{FTN--BC: offline behavior cloning}
\label{sec:bc}
% ===========================================================================

FTN--BC (\code{train\_mt10\_expert.py}) is the purely offline, purely supervised
adaptation. There is no critic, no target network, and no environment interaction
during learning: each task is learned by behavior-cloning a scripted expert from a
fixed dataset of demonstrations. Only the \code{FTNActor} backbone and a single
actor mask $\mask_a$ are used.

\subsection{Data and normalization}
Per task, the MetaWorld scripted expert (from \code{ENV\_POLICY\_MAP}) is rolled
out for \code{n\_demos} episodes; each demo is truncated to
$\texttt{first\_success}+\texttt{demo\_tail}$ frames (full horizon if it never
succeeds), and the $(\obs,\act)$ pairs are concatenated and cached to disk keyed
by the generation parameters so all four conditions reuse identical data. A single
global observation standardizer (mean/std over \emph{all} tasks' data, frozen and
shared) is fit once; actions are already in $[-1,1]$ and left unscaled. The expert
dataset for a task \emph{is} its support set for mask search.

\subsection{BC loss, reconfig, and consolidation}
The reconfig objective is the behavior-cloning loss itself, evaluated on the
(normalized) expert support:
\[
\mathcal L_{\mathrm{BC}}(\masksoft) = \Ehat\Bigl[\,\bigl\|\mu_\phi(\obs;\masksoft) - \act_{\mathrm{expert}}\bigr\|_2^2\,\Bigr],
\]
searched by Algorithm~\ref{alg:ste} and then passed through the lateral smoother
(\S\ref{sec:smooth}). Training is minibatch MSE regression under the committed
mask. The consolidation term (\code{bc\_update}) sums, on top of the current-batch
BC loss, a BC loss over each recalled past chunk \emph{under the mask that chunk
was trained under}:
\[
\mathcal L = \Ehat_{\mathrm{cur}}\bigl[\|\mu_\phi(\obs;\mask_a)-\act\|_2^2\bigr] + \sum_{j}\ \Ehat_{j}\bigl[\|\mu_\phi(\obs;\mask_a^{(j)})-\act^{(j)}\|_2^2\bigr].
\]
Shared sub-networks accumulate gradients from both the new task and the recalled
old tasks---the same consolidation as \S\ref{sec:consolidation}, specialized to
supervised regression. Recalled chunks come from a \code{SupportArchive} that
stores each completed task's (normalized) support and mask and filters by overlap
(Eq.~\eqref{eq:rho}); \code{no\_mask} (overlap $1$) and \code{fixed\_mask}
(overlap $0$) never qualify, so consolidation is automatically a no-op for those
baselines.

\subsection{Per-phase loop and reconfig cadence}
For each task: reset Adam; (\code{fixed\_mask}) assign the disjoint slot; run a
\emph{phase-start} mask reconfig on the expert support and commit; then BC for
\code{bc\_epochs} epochs over shuffled minibatches. Two cadence knobs control
re-search:
\begin{itemize}[leftmargin=1.5em,itemsep=1pt]
  \item \code{reconfig\_train} ($n$ minibatches): in addition to the phase-start
  search, re-search and commit the mask every $n$-th minibatch (warm-started from
  the running logits, on that minibatch); $0$ freezes the mask after phase start.
  \item \code{reconfig\_eval} ($n$ episodes): at eval, the full-FTN condition
  recovers task $t_j$'s mask from its expert support by a cold-started reconfig
  before episode $0$ and again every $n$-th episode ($0$ = once per eval cell,
  reused).
\end{itemize}
After training, the phase archives its support (tagged with the final mask) for
future consolidation, evals every task, and optionally checkpoints the actor
weights, normalizer, and all per-task masks to a single file. Eval under the
default condition recovers each task's mask from its own expert data;
recall/no-mask baselines roll out under the stored/all-ones mask.

% ===========================================================================
\section{Summary}
% ===========================================================================

The three adaptations are one method in three regimes. A fixed-weight bank of
$\Ho$ parallel sub-networks (\S\ref{sec:backbone}) is specialized per task by an
annealed straight-through search for a $\kk$-hot routing mask (\S\ref{sec:ste}),
and past tasks are protected by replaying partially-overlapping archived
experience under its original mask (\S\ref{sec:consolidation}). What varies is the
head and the learning rule attached to the backbone:
\begin{itemize}[leftmargin=1.5em,itemsep=1pt]
  \item \textbf{FTN--DQN} attaches a multi-action $Q$-head, learns by masked TD,
  reconfigures a single mask against returns (default: negative Pearson
  correlation), and replays one overlapping past phase per step under its mask.
  \item \textbf{FTN--DDPG} attaches an actor and a critic, searches a critic mask
  then an actor mask, learns by a summed deterministic-policy-gradient loss over
  the current batch and per-chunk recalled experience (each under its own mask
  pair), and optionally seeds learning with protected expert demonstrations.
  \item \textbf{FTN--BC} attaches only the actor, uses the behavior-cloning MSE as
  both the reconfig objective and the training loss over a fixed offline expert
  dataset, adds lateral spatial smoothing to the searched mask, and consolidates
  via a summed BC loss over overlapping past supports.
\end{itemize}
In every case the mask, not the weights, carries task identity, and the
overlap-filtered replay, not a regularizer, carries cross-task retention.

\end{document}
